%%%%%%%%%%  scan idx 228  |  folio 215  %%%%%%%%%%
induisant sur les faisceaux de sections le morphisme de faisceaux associés à $u_r$.

\subsubsection*{3.2. Définition de l'indice}
\textbf{Théorème 3.2.1.} Soient $S$ un espace localement compact, $p:X\to S$ une variété mixte propre, $E^\bullet$ un complexe elliptique d'ordre $n$ de fibrés vectoriels complexes sur $X$ (1.2.5). Alors le complexe $\mathbb{R}p_*(E^\bullet)$ (qui, d'après (2.5), est isomorphe, à $p_*(E^\bullet)$ dans la catégorie dérivée de celle des $\cO_S\otimes\mathbb C$-Modules) est parfait.

\emph{Preuve.} La question étant locale sur $S$, on peut supposer (2.9) que $X=S\times V$, où $V$ est une variété compacte, et $p=\pr_1$. En outre, on peut supposer (2.10) que l'on a, pour tout $i$, $E^i=\pr_2^*(F^i)$, où $F^i$ est un fibré vectoriel complexe de rang fini sur $V$. D'après (3.1.5), le complexe $p_*(E)$ est associé au complexe de préfaisceaux $\bigl(\underline C'(S,C^\infty(V,F^i)),d_\infty^i\bigr)$ (où les $d^i:E^i\longrightarrow E^{i+1}$ sont les opérateurs différentiels donnés), que nous noterons simplement $\underline C'(S,C^\infty(V,F^\bullet))$. Posons, pour chaque $i\in\Z$,
\[
  m_i=\sum_{j<i} n_j
  \qquad (\text{on convient que } n_i=0 \text{ si } E^i=0).
\]
Pour chaque $r\in\Z$, on a (3.1.5) un complexe de préfaisceaux $\bigl(\underline C'(S,C_2^{r-m_i}(V,F^i)),d_{r-m_i}^i\bigr)$, que nous noterons $\underline C'(S,C_2^r(V,F^\bullet))$.
Les $\underline C'(S,C_2^r(V,F^\bullet))$ forment un système projectif, et l'on a (3.1.5)~:
\begin{equation}\tag{3.2.1.1}
  \underline C'(S,C^\infty(V,F^\bullet))=\varprojlim_r \underline C'(S,C_2^r(V,F^\bullet)).
\end{equation}
Enfin, nous noterons $F_r^\bullet$ le complexe de fibrés hilbertiens triviaux sur $S$ défini par $\varphi_{r-m_i}d^i$, et $\underline F_r^\bullet$ le complexe des faisceaux de sections correspondant, qui n'est autre (3.1.5) que le complexe de faisceaux associés à $\underline C'(S,C_2^r(V,F^\bullet))$. Cela posé, pour démontrer le théorème, il suffit de prouver les assertions suivantes~:

(a) Pour $r\in\Z$, le morphisme canonique
